% Discussion D1 — which property of a foundation model predicts patient-state representation? (synced to P0 round-17 state, 2026-09-09)
\subsection{Which model property tracks patient-state representation?}
\label{ssec:discussion_property}

Architecture family is not identifiable in these data: although the
within-family spread among bidirectional transformer encoders (from
\model{Geneformer} 0.786 to \model{UCE} 0.661) exceeds the gap to the only
generative model, every between-family contrast is small against the roughly
0.10-AUC resolution of a 25-event benchmark, and the strongest between-family
gaps (\model{STATE}$-$\model{scGPT} $+0.055$;
\model{Geneformer}$-$\model{scGPT} $+0.068$; exact DeLong $p = 0.014$ and
$0.008$, in agreement with the paired bootstrap;
Supplementary Note~\ref{sn:se_check}) sit exactly at that resolution's
edge. We therefore do not claim architecture-level attribution.

The one confirmatory separation, the two leaders versus \model{UCE}
(range-statistic $P = 0.011$; Supplementary Note~\ref{sn:range_test}), is a
statement about the losing model before it is a statement about the winning
ones. \model{UCE} is the only method that never significantly outscored the
count-plus-disease null model ($p = 0.157$;
Supplementary Note~\ref{sn:method_vs_null}), so the content of the gap is
that \model{UCE} is unsuited to this task, not that the leaders carry a
pretraining advantage, a claim the primary test rejects
(Section~\ref{ssec:r1_within_covid}). The natural attribution hypothesis for
such a gap, blood and immune representation in the pretraining corpus, was
measured directly and is falsified on both ends
(Supplementary Note~\ref{sn:corpus_platelet}). The three models whose
platelet content can be numerically bounded against public corpus records
(\model{scPRINT-2} 0.120--0.125\%, \model{scGPT} 0.125\%, \model{UCE}
0.120\%) share essentially identical fractions yet span 0.079 patient-level
AUC, 63\% of the entire between-model spread. Of the two leaders,
\model{STATE}'s corpus status is itself bounded and is the lowest documented
of all --- about two-thirds of its corpus (the scBaseCount and Tahoe
components) is platelet-free at both the annotation and quality-control
layers, and the remainder shares the Census bound, capping its platelet
content below 0.05\% (Table~\ref{tab:models}) --- while \model{Geneformer}
is the one model with no cell-type-level corpus composition in any published
material. A leader at a documented near-zero and a laggard at 0.120\%
produce the same ordering as one leader with a documented corpus and one
without: wherever corpus platelet content can actually be checked, it does
not track performance.

Three further facts bound any property claim. First, corpus metrics must be
read in one unit, and none exists across all five models: the fraction of
pretraining cells from blood and bone marrow is documented as dataset counts
for \model{Geneformer} (49 of 561 source datasets, under a 25\%-per-tissue
cap), as a sample fraction for \model{UCE} (6.6\%), and not at all for the
rest (Supplementary Table~\ref{tab:corpus_fractions}). Second,
\model{scPRINT-2} contradicts any simple identity-objective story: its
objective includes explicit cell-label prediction, the attribute assigned to
the losing model, yet it ranks third, above the baseline on point estimates;
and the from-scratch PCA-200 baseline, with no pretraining at all, lands
mid-table and is statistically inseparable from two of the five models.
Third, the ordering does not persist on the in-distribution monocyte control:
\model{scGPT} moves from fourth on platelets to second on monocytes
(Section~\ref{ssec:r4_monocyte}), so no corpus-level constant can explain the
ranking on both cell types.

We therefore state the property question as a set of benchmark-testable
hypotheses, one of which now carries a recorded falsification. Supplementary
Table~\ref{tab:property_multiplicity} enumerates the candidate attributes we
checked (corpus blood/marrow fraction, corpus disease-condition fraction,
objective loss target, tokenization, corpus size, parameter count, embedding
dimension, vocabulary size, species coverage, and release year) and reports
how many separate the observed ranking equally well; the corpus-composition
hypothesis above remains in the set as an untested candidate, now known to be
falsified where it is measurable and untestable where it would matter.
Resolving any such property at the 0.03--0.05 AUC scale where attribution
lives requires on the order of $10^3$ patients at this event rate; the
present benchmark is designed so that any new model or new candidate property
can be inserted and tested under the identical protocol.
